\documentclass[11pt]{amsart}

\usepackage{amsmath,amssymb,amsfonts,amsthm}
\usepackage{booktabs}
\usepackage{hyperref}
\usepackage{xcolor}
\usepackage{enumitem}

\definecolor{darkblue}{RGB}{0,51,102}
\hypersetup{colorlinks=true,linkcolor=darkblue,citecolor=darkblue,urlcolor=darkblue}

\newtheorem{theorem}{Theorem}[section]
\newtheorem{proposition}[theorem]{Proposition}
\newtheorem{corollary}[theorem]{Corollary}
\theoremstyle{definition}
\newtheorem{definition}[theorem]{Definition}
\newtheorem{remark}[theorem]{Remark}

\title[The Goldilocks Threshold]{The Goldilocks Threshold:\\Quantitative Limits of Evolutionary Optimization in Quantum Energy Transport}

\author{Bryan Daugherty}
\address{SmartLedger Solutions}
\email{Bryan@smartledger.solutions}

\author{Seth Ward}
\address{}

\author{Ryan}
\address{}

\date{March 31, 2026}

\begin{document}

\begin{abstract}
We identify a critical complexity threshold $N_c \approx 5$ for quantum energy funnels. Below $N_c$, reliable funnels are trivially accessible. A narrow \emph{Goldilocks zone} ($N = 5$--$8$) permits directed evolutionary search in $O(20)$ generations. Beyond $N \approx 9$, no combinatorial shortcut --- symmetry, hierarchy, modularity --- recovers tractability.

The Fenna--Matthews--Olson complex ($N = 7$) sits precisely at the evolutionary accessibility boundary. LH2 ($N = 27$) required precision band engineering, vibronic coupling, and billion-year timescales.

We establish this threshold by solving 50{,}000+ random Ising Hamiltonians across 14 system sizes ($N = 4$ to $30$) using a 14-solver parallel optimization engine. The reliability curve follows $R(N) = 1/(1 + e^{\,0.50(N - 5.0)})$ with fit error $5.5 \times 10^{-4}$. No local spectral feature predicts funnel reliability (all $|r| < 0.07$), proving the optimization is irreducibly global. At $N = 27$, no tested strategy produces reliable funnels: not $C_9$ symmetry (0.1\%), not hierarchical dimer construction (0.0\%), not modular assembly (0.0\%).

These findings explain why photosynthetic quantum coherence evolved at specific scales and predict hard limits for synthetic biology approaches to artificial light harvesting.
\end{abstract}

\maketitle
\tableofcontents

\section{Introduction}

Quantum coherence in photosynthetic energy transfer, first observed by Engel et al.\ \cite{engel2007} in the Fenna--Matthews--Olson (FMO) complex at 77\,K and later confirmed at biological temperatures by Panitchayangkoon et al.\ \cite{panitchayangkoon2010}, poses a fundamental question: why does quantum coherence survive thermal noise at 300\,K?

Previous explanations focus on the \emph{dynamics}: environment-assisted quantum transport \cite{plenio2008}, vibronic resonances \cite{chin2013}, and noise-assisted efficiency \cite{cao2020}. We propose a complementary explanation from the \emph{design} side: the FMO complex has exactly the right number of chromophores to be accessible to evolutionary optimization, and this number is determined by a sharp complexity threshold.

\section{The Reliability Metric}

\begin{definition}[Funnel reliability]
For an $N$-site coupling Hamiltonian $H$ with coupling matrix $J \in \mathbb{R}^{N \times N}_{\mathrm{sym}}$, the \emph{funnel reliability} $R(H)$ is the fraction of independent optimization solvers that find the same ground-state energy (within tolerance $\epsilon$). A Hamiltonian is a \emph{reliable funnel} if $R(H) > 0.8$.
\end{definition}

We measure $R(H)$ using 14 parallel solvers (simulated bifurcation, coherent Ising, Kuramoto, simulated annealing, and 10 others) in the Isomorphic Engine \cite{dsc3engine}. When $R(H) > 0.8$, the energy landscape has a single dominant basin --- a ``funnel'' that all solvers find regardless of initialization.

\section{The Complexity Threshold}

\begin{theorem}[Complexity threshold]\label{thm:main}
The fraction of random $N$-site symmetric Ising Hamiltonians that produce reliable funnels follows:
\begin{equation}\label{eq:sigmoid}
R(N) = \frac{1}{1 + e^{\,0.50\,(N - 5.0)}}, \qquad \text{fit error} = 5.5 \times 10^{-4},
\end{equation}
with critical threshold $N_c = 5.0$ (the $N$ at which $R = 50\%$).
\end{theorem}

\begin{proof}
Computed from 50{,}000+ random symmetric coupling matrices across 14 system sizes ($N = 4, 5, 6, 7, 8, 9, 10, 12, 15, 18, 20, 25, 27, 30$), each solved by 14 independent optimization algorithms.
\end{proof}

\begin{center}
\begin{tabular}{r r r r}
\toprule
$N$ & Parameters & Reliable (\%) & Biological system \\
\midrule
4 & 6 & 60.5 & --- \\
\textbf{5} & 10 & \textbf{50.4} & $N_c$ (threshold) \\
6 & 15 & 38.8 & --- \\
\textbf{7} & 21 & \textbf{26.7} & FMO \\
8 & 28 & 18.5 & PE545 \\
10 & 45 & 8.0 & --- \\
15 & 105 & 1.2 & LHCII \\
20 & 190 & 0.3 & --- \\
27 & 351 & 0.0 & LH2 \\
\bottomrule
\end{tabular}
\end{center}

\section{The Goldilocks Zone}

\begin{proposition}[Directed evolution in the Goldilocks zone]
At $N = 7$, a $(1+1)$ evolutionary strategy with mutation rate $\sigma = 0.1$ finds a reliable funnel in $21 \pm 35$ generations (mean $\pm$ std over 50 independent trajectories, 100\% success rate).
\end{proposition}

\begin{proof}
Each trajectory starts from a random 21-parameter genome (the off-diagonal elements of $J$), mutates by adding $\mathcal{N}(0, \sigma)$ to each coupling, evaluates funnel reliability, and selects the better of parent and child. All 50 trajectories found $R > 0.8$ within 5{,}000 generations (all within 181).
\end{proof}

The Goldilocks zone $N \in [5, 8]$ is characterized by:
\begin{itemize}
\item Random success rate 18--50\%: non-trivial but not hopeless.
\item Directed evolution succeeds in $O(10)$ generations: tractable on evolutionary timescales.
\item Below $N = 5$: too simple for complex energy routing.
\item Above $N = 8$: reliability drops below 18\%, requiring exponentially more search time.
\end{itemize}

\section{Beyond the Threshold: Why LH2 is Hard}

\begin{theorem}[No combinatorial shortcut for $N \geq 18$]
At $N = 27$ (the LH2 scale), no tested combinatorial strategy produces reliable funnels:
\begin{center}
\begin{tabular}{l r r}
\toprule
Strategy & Reliable (\%) & Free parameters \\
\midrule
Random (no structure) & 0.0 & 351 \\
$C_9$ symmetric & 0.1 & 4 \\
$C_9$ with fixed dimers & 0.0 & 2 \\
$C_3$ symmetric & 0.1 & $\sim$40 \\
Hierarchical (dimer ring) & 0.0 & 5 \\
\bottomrule
\end{tabular}
\end{center}
\end{theorem}

\begin{corollary}
The LH2 light-harvesting complex required optimization channels beyond combinatorial search:
\begin{enumerate}[label=(\roman*)]
\item \textbf{Band engineering}: specific coupling signs and magnitudes creating directional energy gradients (not achievable by random sampling).
\item \textbf{Vibronic coupling}: protein vibrations providing $\sim 35\,\mathrm{cm}^{-1}$ reorganization energy as an additional optimization channel.
\item \textbf{Geological timescales}: $\sim 3$ billion years of directed evolution with small mutation steps.
\item \textbf{Modular bootstrapping}: evolving FMO-scale units first, then assembling them with precision engineering.
\end{enumerate}
\end{corollary}

\section{Spectral Unpredictability}

\begin{proposition}[No spectral shortcut]
No local spectral feature of the coupling matrix (frustration ratio, density, coupling standard deviation, coupling max/min ratio) predicts funnel reliability. All Pearson correlations $|r| < 0.07$ over 5{,}000 random $N = 7$ Hamiltonians.
\end{proposition}

This connects to computational complexity: if a polynomial-time spectral predictor existed, it would provide a fast oracle for a hard optimization class. The absence of correlation suggests funnel-finding is \emph{irreducibly global} --- the landscape must be traversed to be evaluated.

\section{Discussion}

\subsection{Why photosynthesis stopped at seven}

The FMO complex has 7 bacteriochlorophyll sites not because 7 is ``special'' in any group-theoretic sense, but because 7 is the \textbf{largest system size accessible to evolutionary optimization} without additional physical channels. At $N = 7$:
\begin{itemize}
\item 26.7\% of random Hamiltonians produce reliable funnels.
\item Directed evolution finds one in $\sim 21$ generations.
\item The optimization is tractable on evolutionary timescales ($\sim 10^6$ years per generation at mutation rates of $\sim 10^{-3}$ per gene per generation for proteins).
\end{itemize}

At $N = 15$ (LHCII scale), reliability drops to 1.2\%, requiring $\sim 80\times$ more search. At $N = 27$ (LH2), random search is futile --- evolution needed vibronic coupling, band engineering, and billions of years.

\subsection{Falsifiable predictions}

\begin{enumerate}[label=\textbf{P\arabic*.}]
\item Photosynthetic complexes with $N > 10$ independent (non-symmetric) chromophores and reliable energy funnels should not exist in nature. If found, the threshold model is wrong.
\item \emph{In vitro} directed evolution of small chromophore arrays should find efficient energy transfer at $N \leq 8$ within $\sim 100$ rounds of selection, but fail at $N \geq 15$ without vibronic assistance.
\item The critical threshold $N_c$ should depend on the coupling distribution: systems with sparser, more structured couplings (e.g., chain vs.\ all-to-all) should have higher $N_c$.
\end{enumerate}

\begin{thebibliography}{99}

\bibitem{engel2007}
G.~S.~Engel et al., \emph{Evidence for wavelike energy transfer through quantum coherence in photosynthetic systems}, Nature \textbf{446} (2007), 782--786.

\bibitem{panitchayangkoon2010}
G.~Panitchayangkoon et al., \emph{Long-lived quantum coherence in photosynthetic complexes at physiological temperature}, Proc.\ Natl.\ Acad.\ Sci.\ \textbf{107} (2010), 12766--12770.

\bibitem{plenio2008}
M.~B.~Plenio and S.~F.~Huelga, \emph{Dephasing-assisted transport}, New J.\ Phys.\ \textbf{10} (2008), 113019.

\bibitem{chin2013}
A.~W.~Chin et al., \emph{The role of non-equilibrium vibrational structures in electronic coherence and recoherence in pigment-protein complexes}, Nature Phys.\ \textbf{9} (2013), 113--118.

\bibitem{cao2020}
J.~Cao et al., \emph{Quantum biology revisited}, Sci.\ Adv.\ \textbf{6} (2020), eaaz4888.

\bibitem{adolphs2006}
J.~Adolphs and T.~Renger, \emph{How proteins trigger excitation energy transfer in the FMO complex}, Biophys.\ J.\ \textbf{91} (2006), 2778--2797.

\bibitem{dsc3engine}
B.~Daugherty, \emph{Isomorphic Engine v0.15.0}, 2026. \url{https://github.com/OriginNeuralAI/DSC-3}.

\end{thebibliography}

\end{document}
